\section{Conclusion}\label{sec:conclusion}

CinePal contributes a complete, reproducible system for clustering a catalogue by conversation,
where an oracle's acceptance, not any intrinsic metric, is the objective. Its technical
core is a multi-agent architecture in which a deterministic coordinator routes each turn through
narrow, read-only LLM agents and a content-addressed snapshot tree that makes every session
replayable, fed by a multi-modal embedding pipeline and exercised through an automated
oracle-and-judge evaluation harness. Measured against a single-prompt baseline, the conversational
system clearly chooses more appropriate operations and follows the oracle's intent more faithfully,
at no extra cost per turn, while the baseline's apparent tidiness is a side effect of
under-segmentation. The system's weaknesses are concrete and local, namely wrong-film explanations,
bloated labels, and unhonoured split counts, which makes them actionable rather than
fundamental. The build, the trace it leaves behind, and the findings it makes legible are the
contribution.
